\documentclass[12pt,a4paper]{article}

% --- 宏包配置 ---
\usepackage[UTF8]{ctex}
\usepackage{amsmath, amssymb, amsfonts, amsthm}
\usepackage{bm}       % 提供 \bm 命令用于矢量加粗
\usepackage{geometry}    % 页边距
\usepackage{enumitem}    % 列表格式
\usepackage{tcolorbox}   % 重点框选
\usepackage{booktabs}    % 表格优化
\usepackage{longtable}   % 长表格
\usepackage{titlesec}    % 标题定制
\usepackage[
    colorlinks=true,
    linkcolor=blue,
    unicode=true,          % 关键：使用 unicode 编码处理书签
    pdfpagemode=UseNone,   % 避免打开时自动弹出书签栏导致某些预览器崩溃
    pdfstartview=FitH
]{hyperref}
\setcounter{secnumdepth}{4}
\setcounter{tocdepth}{4}
\geometry{left=2cm, right=2cm, top=2.5cm, bottom=2.5cm}

% --- 自定义矢量命令 ---
\newcommand{\vect}[1]{\mathbf{#1}}                   % 矢量加粗
\newcommand{\unitv}[1]{\hat{\mathbf{#1}}}            % 单位矢量
\newcommand{\diff}{\mathrm{d}}                       % 直体微分符号 d
\newcommand{\der}[2]{\frac{\mathrm{d} #1}{\mathrm{d} #2}}  % 全导数
\newcommand{\pder}[2]{\frac{\partial #1}{\partial #2}} % 偏导数
\newcommand{\dder}[2]{\frac{\mathrm{d}^2 #1}{\mathrm{d} #2^2}}   % 二阶全导数

\title{\textbf{矢量力学}}
\author{魏舟}
\date{2026年1月}

\begin{document}

\maketitle
本文档根据我自己的笔记和Gemini的矢量力学课程资料整理而成，涵盖静力学与运动学两大部分内容，旨在为学习和复习矢量力学提供系统化的参考资料。
\newpage

\tableofcontents

\newpage
% ==============================================================================
% 第一章：静力学公理与物体受力分析
% ==============================================================================
\section{静力学公理及其力学本质}

静力学研究物体在力系作用下的平衡规律。静力学逻辑建立在以下五条公理之上：

\subsection{静力学基本公理}
\begin{itemize}[leftmargin=2em]
    \item \textbf{公理一：二力平衡原理}。作用在刚体上的两个力，使刚体保持平衡的充要条件是：两力大小相等、方向相反、作用在同一直线上。
    \item \textbf{公理二：加减平衡力系原理}。在已知力系上加上或减去任意平衡力系，不改变原力系对刚体的效应。
    \begin{itemize}
        \item \textit{推论（力的可传性）：} 作用于刚体上的力可沿其作用线移动到刚体内任意点。
    \end{itemize}
    \item \textbf{公理三：力的平行四边形法则}。作用于物体上同一点的两个力，可以合成为一个合力。合力的矢量等于两个分力的矢量和：$\vect{F}_R = \vect{F}_1 + \vect{F}_2$。
    \item \textbf{公理四：作用与反作用定律}。两物体间的相互作用力，总是大小相等、方向相反、作用在同一直线上。
    \item \textbf{公理五：刚化原理}。若变形体在已知力系作用下平衡，则将其“刚化”为刚体后，平衡状态保持不变。
\end{itemize}

\subsection{受力分析标准化步骤 (Standardized Procedure)}
在处理复杂工程结构时，必须严格执行以下受力分析流程，严禁凭直觉画力。

\begin{tcolorbox}[colback=blue!5!white, colframe=blue!75!black, title=标准化四步法]
\begin{enumerate}
    \item \textbf{确定研究对象}：根据待求量，选取单个物体或物体系（整体）。
    \item \textbf{取隔离体}：将研究对象从周围约束中脱离，画出独立的轮廓图。
    \item \textbf{画主动力}：画出重力（$G=mg$）、已知载荷、流体压力等。
    \item \textbf{画约束力}：根据约束类型，逐一分析并画出约束力。
\end{enumerate}
\end{tcolorbox}

\subsection{典型约束及其矢量特性}

\begin{longtable}{lll}
    \toprule
    \textbf{约束类型} & \textbf{约束力特性} & \textbf{数学分量表示} \\
    \midrule
    光滑面约束 & 沿法线方向，指向受力体 & $F_n$ \\
    柔体约束 (绳索) & 沿中心线，背离受力体 & $F_T$ \\
    径向轴承 (铰链) & 通过销钉中心，方向未知 & $F_x, F_y$ \\
    球铰链 (空间) & 通过球心，空间任一方向 & $F_x, F_y, F_z$ \\
    固定支座 & 限制所有位移与转动 & $F_x, F_y, M$ \\
    二力杆 & 沿两端铰链连线 & $F_{AB}$ \\
    \bottomrule
\end{longtable}

% ==============================================================================
% 第二章：力系的简化
% ==============================================================================
\section{力系的简化与等效转化}

\subsection{力对点之矩 (Moment of Force about a Point)}
力 $\vect{F}$ 对点 $O$ 之矩是一个定位矢量：
\begin{equation}
    \vect{M}_O(\vect{F}) = \vect{r} \times \vect{F}
\end{equation}
其三维分量形式（行列式展开）：
\begin{equation}
    \vect{M}_O(\vect{F}) = 
    \begin{vmatrix}
        \unitv{i} & \unitv{j} & \unitv{k} \\
        x & y & z \\
        F_x & F_y & F_z
    \end{vmatrix} = (yF_z - zF_y)\unitv{i} + (zF_x - xF_z)\unitv{j} + (xF_y - yF_x)\unitv{k}
\end{equation}

\subsection{力对轴之矩 (Moment of Force about an Axis)}
力对轴之矩是力对该轴上任一点之矩在轴上的投影：
\begin{equation}
    M_z(\vect{F}) = [\vect{M}_O(\vect{F})] \cdot \unitv{k}
\end{equation}
\textbf{性质：}若力 $\vect{F}$ 与轴平行或力 $\vect{F}$ 的作用线与轴相交，则力对该轴之矩为零。

\subsection{力偶 (Couple)}
力偶是由两个等大、反向、不共线的平行力组成的力系。
\begin{itemize}
    \item \textbf{力偶矩矢量}：$\vect{M} = \vect{r}_{BA} \times \vect{F}$。
    \item \textbf{特性}：力偶矩与简化中心的选择无关，是一个自由矢量。
    \item \textbf{等效性}：只要力偶矩矢量相同，力偶对刚体的效应就完全相同。
\end{itemize}

\subsection{力系向一点简化的普遍结论}
任意空间力系向一点 $O$ 简化的结果为一个主矢 $\vect{R}'$ 和一个主矩 $\vect{M}_O$：
\begin{align}
    \vect{R}' &= \sum \vect{F}_i \\
    \vect{M}_O &= \sum \vect{M}_O(\vect{F}_i) = \sum (\vect{r}_i \times \vect{F}_i)
\end{align}

\begin{tcolorbox}[colback=red!5!white, colframe=red!75!black, title=主矢与主矩的区别]
主矢 $\vect{R}'$ 是力系的矢量和，与简化中心的选择\textbf{无关}；\\
主矩 $\vect{M}_O$ 与简化中心的选择\textbf{有关}，满足换点公式：
$\vect{M}_{O'} = \vect{M}_O + \vect{r}_{O'O} \times \vect{R}'$。
\end{tcolorbox}

% ==============================================================================
% 第三章：力系的平衡
% ==============================================================================
\section{力系的平衡条件与平衡方程}

\subsection{空间一般力系的平衡方程}
空间力系平衡的充要条件是：主矢和对任一点的主矩均为零。
\begin{align}
    \sum F_x = 0, \quad \sum F_y = 0, \quad \sum F_z = 0 \\
    \sum M_x = 0, \quad \sum M_y = 0, \quad \sum M_z = 0
\end{align}

\subsection{平面一般力系的平衡方程}
对于平面力系，常见的平衡方程形式有三种：
\begin{enumerate}
    \item \textbf{基本形式}（一力矩式）：
    $\sum F_x = 0, \sum F_y = 0, \sum M_O = 0$。
    \item \textbf{二力矩式}：
    $\sum F_x = 0, \sum M_A = 0, \sum M_B = 0$。（限制：$AB$ 连线不垂直于 $x$ 轴）
    \item \textbf{三力矩式}：
    $\sum M_A = 0, \sum M_B = 0, \sum M_C = 0$。（限制：$A, B, C$ 不共线）
\end{enumerate}

\subsection{考虑摩擦的平衡问题 (Static Friction)}

\begin{itemize}
    \item \textbf{库仑摩擦定律}：$F_f \leq F_{s,max} = f_s \cdot F_n$。
    \item \textbf{全反力与摩擦角}：当摩擦力达到最大值时，全反力 $\vect{F}_R$ 与法线的夹角称为摩擦角 $\phi_s$，满足 $\tan \phi_s = f_s$。
    \item \textbf{自锁条件}：当主动合力的作用线在摩擦角之内时，无论力多大，物体始终保持静止。
\end{itemize}

\subsection{空间平衡计算详解案例}
设一长方体框架，在顶点 $C(a, b, c)$ 处受力 $\vect{F} = F_x \unitv{i} + F_y \unitv{j} + F_z \unitv{k}$。求对原点 $O$ 的力矩。
\begin{align}
    \vect{M}_O(\vect{F}) &= 
    \begin{vmatrix}
        \unitv{i} & \unitv{j} & \unitv{k} \\
        a & b & c \\
        F_x & F_y & F_z
    \end{vmatrix} \notag \\
    &= (bF_z - cF_y)\unitv{i} + (cF_x - aF_z)\unitv{j} + (aF_y - bF_x)\unitv{k}
\end{align}
若该力系平衡，则上述三个分量必须分别抵消掉支座产生的约束力矩。

% ==============================================================================
% 第四章：静力学经典题型分析方法
% ==============================================================================
\section{静力学经典题型分析方法}

\subsection{题型一：物体系的平衡问题（拆解与整体法）}
对于由多个杆件、滑轮或几何体组成的系统，核心在于**灵活切换研究对象**。

\textbf{分析方法：}
\begin{enumerate}
    \item \textbf{整体法}：先将整个系统作为研究对象，求出外部支座的约束力。
    \item \textbf{拆解法}：在铰链、接触点处将物体拆开。注意应用“作用与反作用定律”，在拆开的接触面上，两物体受到的力必须**等大、反向、共线**。
    \item \textbf{约束性质判定}：优先寻找“二力杆”。
\end{enumerate}



\subsection{题型二：考虑滑动摩擦的临界平衡问题}
这类题目的特点是：物体处于“将动未动”的状态，要求确定力或角度的取值范围。

\textbf{分析方法：}
\begin{enumerate}
    \item \textbf{建立平衡方程}：按照常规静力学方程列式。
    \item \textbf{引入摩擦补充方程}：在临界状态，补充 $F_f = f_s F_N$。
    \item \textbf{几何法（摩擦角法）}：对于单体受力，利用全反力 $\vect{F}_R$ 必须落在摩擦锥（Friction Cone）内的性质，通过正弦定理求解。
\end{enumerate}



\subsection{题型三：桁架结构的内力计算（节点法与截面法）}
哈工大教材中重点考察的是简单桁架，其假设所有杆件均为二力杆。

\textbf{分析方法：}
\begin{itemize}
    \item \textbf{节点法 (Method of Joints)}：
    取单个节点为研究对象，利用 $\sum F_x = 0, \sum F_y = 0$ 求解。适用于求所有杆件内力。
    \item \textbf{截面法 (Method of Sections)}：
    用虚构截面将桁架切开，取一部分为研究对象。利用三个平衡方程求解被切断的杆件内力（通常不超过3根）。
    \textit{技巧：求某杆内力时，将矩心选在另外两根未知力杆的交点上。}
\end{itemize}

\subsection{题型四：空间力系的简化与平衡推导}
这是矢量力学基本功的体现，通常涉及空间支架或轴系传动。

\textbf{分析方法：}
\begin{enumerate}
    \item \textbf{矢量坐标化}：将所有力 $\vect{F}_i$ 表示为坐标分量形式 $\vect{F}_i = F_{ix}\unitv{i} + F_{iy}\unitv{j} + F_{iz}\unitv{k}$。
    \item \textbf{力矩投影}：
    对于复杂的空间结构，直接用行列式求矩：
    \begin{equation}
        \vect{M}_O = \sum (\vect{r}_i \times \vect{F}_i)
    \end{equation}
    \item \textbf{方程组联立}：得到六个代数方程。若支座约束超过六个，则属于静不定问题（理论力学阶段通常不涉及，或需结合对称性处理）。
\end{enumerate}



\section{静力学部分的总结与辨析}

\begin{tcolorbox}[colback=green!5!white, colframe=green!50!black, title=解题避坑指南]
    \begin{itemize}
        \item \textbf{力偶只能出力偶去平衡}：力偶不能与一个力平衡。如果系统受一个力偶作用，必然存在另一个力偶或一组产生等效力偶矩的约束力。
        \item \textbf{二力杆的判定}：二力杆不一定是直杆，只要只在两点受力且不计自重，内力必沿两点连线。
        \item \textbf{三力平衡汇交}：若物体受三个力平衡，且其中两个力相交，则第三个力必过交点。利用此几何特性可大幅简化方程。
    \end{itemize}
\end{tcolorbox}


% ==============================================================================
% 第二部分：运动学 (Kinematics)
% ==============================================================================
\section{运动学基础与矢量导数}
运动学研究物体的运动几何性质，而不考虑引起运动的力。在矢量力学中，其核心工具是矢量对时间的导数。

\subsection{自然坐标系下的运动描述}
当点在空间曲线运动时，引入切向单位矢量 $\unitv{\tau}$、法向单位矢量 $\unitv{n}$ 和副法向单位矢量 $\unitv{b}$。
\begin{itemize}
    \item \textbf{速度矢量}：$\vect{v} = v \unitv{\tau}$
    \item \textbf{加速度分解}：
    \begin{equation}
    \vect{a} = \der{\vect{v}}{t} = \der{v}{t}\unitv{\tau} + v\der{\unitv{\tau}}{t} = \vect{a}_t + \vect{a}_n
\end{equation}
    其中，切向加速度 $a_t = \dot{v}$ 描述速度大小的变化；法向加速度 $a_n = \frac{v^2}{\rho}$ 描述速度方向的变化，$\rho$ 为轨迹曲率半径。
\end{itemize}

% ==============================================================================
\section{点的合成运动 (Kinematics of Particles in Relative Motion)}

这是理论力学的重点。我们需要在定参考系（$Oxyz$）和动参考系（$O'x'y'z'$）之间建立联系。

\subsection{动点与动系选择的四条准则}
为了使物理图像清晰且计算可行，选择动点和动系时应遵循：
\begin{enumerate}
    \item \textbf{准则一：相对运动的可观察性}。动点相对于动系的轨迹必须是已知或易于确定的几何线（如直线槽、圆弧轨道）。
    \item \textbf{准则二：动系固连性}。动系必须牢固地连在某个具有已知运动特性的刚体上，通常该刚体作定轴转动或平动。
    \item \textbf{准则三：独立性}。动点不能选择动系所在的刚体上的点。
    \item \textbf{准则四：唯一性}。动点在特定瞬时的绝对位置应当是物理上唯一的。
\end{enumerate}

\subsection{速度合成定理}
设 $\vect{r}_a$ 为动点在定系中的位置，$\vect{r}_O'$ 为动系原点在定系中的位置，$\vect{r}_r$ 为动点在动系中的位置。则：
$\vect{r}_a = \vect{r}_O' + \vect{r}_r$。
对时间求导得：
\begin{equation}
    \vect{v}_a = \vect{v}_e + \vect{v}_r
\end{equation}
其中 $\vect{v}_e$ 为牵连速度（动系上与动点重合的牵连点之速度），$\vect{v}_r$ 为相对速度。

\subsection{科氏加速度 $\vect{a}_k$ 的深度推导}

当牵连运动包含转动（$\vect{\omega}_e \neq 0$）时，加速度合成定理需包含科氏项。

\textbf{推导过程：矢量导数变换}

设动系角速度为 $\vect{\omega}_e$。对于任何动系中的矢量 $\vect{A}$，其在定系中的导数为：
\begin{equation}
    \left( \der{\vect{A}}{t} \right)_{fix} = \left( \der{\vect{A}}{t} \right)_{mov} + \vect{\omega}_e \times \vect{A}
\end{equation}
将此算子应用于速度合成式 $\vect{v}_a = \vect{v}_e + \vect{v}_r = \vect{v}_O' + \vect{\omega}_e \times \vect{r}_r + \vect{v}_r$：
\begin{align}
    \vect{a}_a &= \der{\vect{v}_O'}{t} + \der{\vect{\omega}_e}{t} \times \vect{r}_r + \vect{\omega}_e \times \der{\vect{r}_r}{t} + \der{\vect{v}_r}{t} \notag \\
    &= \vect{a}_O' + \vect{\alpha}_e \times \vect{r}_r + \vect{\omega}_e \times (\vect{v}_r + \vect{\omega}_e \times \vect{r}_r) + (\vect{a}_r + \vect{\omega}_e \times \vect{v}_r) \notag \\
    &= \underbrace{\vect{a}_O' + \vect{\alpha}_e \times \vect{r}_r + \vect{\omega}_e \times (\vect{\omega}_e \times \vect{r}_r)}_{\vect{a}_e} + \vect{a}_r + \underbrace{2\vect{\omega}_e \times \vect{v}_r}_{\vect{a}_k}
\end{align}

\textbf{科氏加速度的物理意义与消失条件}

$\vect{a}_k$ 的存在源于两点：1. 相对位置改变引起的牵连速度改变；2. 牵连转动引起的相对速度方向改变。
\begin{itemize}
    \item \textbf{消失条件}：$\vect{\omega}_e = 0$（牵连平动）；$\vect{v}_r = 0$（动点相对静止）；$\vect{\omega}_e \parallel \vect{v}_r$（沿转轴滑动）。
\end{itemize}

% ==============================================================================
\section{刚体平面运动 (Kinematics of Planar Motion)}

平面运动是指刚体上各点到某一固定平面的距离始终保持不变。

\subsection{速度分析：基点法与瞬心法}
\subsubsection{基点法 (Base Point Method)}
选取刚体上一点 $B$ 为基点，点 $A$ 的速度为：
\begin{equation}
    \vect{v}_A = \vect{v}_B + \vect{v}_{AB} = \vect{v}_B + \vect{\omega} \times \vect{r}_{BA}
\end{equation}
此式揭示了平面运动可分解为：随基点的平动 + 绕基点的转动。

\subsubsection{速度瞬心法 (Instantaneous Center of Rotation)}

在任一瞬时，若刚体角速度 $\vect{\omega} \neq 0$，必存在唯一的点 $P$（瞬心），其速度 $\vect{v}_P = 0$。
此时，刚体上任意点 $i$ 的速度模长为 $v_i = \omega \cdot \rho_{Pi}$，方向垂直于 $Pi$ 连线。
\subsubsection{寻找速度瞬心 (ICR) 的几何与矢量方法}

速度瞬心是在某一瞬时，刚体上或其延伸体上速度为零的点 $P$。判定瞬心位置是计算刚体角速度 $\omega$ 的前提。

\paragraph{方法一：已知两点速度方向 (垂直线交点法)}
这是最常用的方法。若已知刚体上 $A, B$ 两点的速度 $\vect{v}_A$ 和 $\vect{v}_B$ 的方向：
\begin{enumerate}
    \item 分别过 $A$ 点和 $B$ 点作速度矢量 $\vect{v}_A$ 和 $\vect{v}_B$ 的垂线。
    \item 两垂线的交点即为该瞬时的速度瞬心 $P$。
\end{enumerate}
\textit{注意：若两垂线平行且不重合，则说明刚体此时角速度 $\omega = 0$，处于瞬时平动状态。}



\paragraph{方法二：已知两点速度大小且方向平行 (相似三角形法)}
当 $A, B$ 两点的速度 $\vect{v}_A \parallel \vect{v}_B$ 且均垂直于 $AB$ 连线时：
\begin{itemize}
    \item \textbf{同向情况}：连接 $\vect{v}_A$ 和 $\vect{v}_B$ 的末端，交 $AB$ 连线（或延长线）于点 $P$。
    \item \textbf{反向情况}：连接两速度矢量末端，交点 $P$ 必在 $AB$ 段之间。
\end{itemize}
利用相似三角形关系求位置：$\frac{v_A}{AP} = \frac{v_B}{BP} = \omega$。



\paragraph{方法三：纯滚动约束法 (接触点法)}
这是解决车轮、齿轮等问题的捷径。
\begin{itemize}
    \item 若一圆盘在地面上作**纯滚动**（无滑动），则圆盘与地面的接触点 $C$ 瞬时速度为零，即 $P = C$。
    \item \textbf{推广}：若两个刚体在接触点处无相对滑动，则接触点是它们相对运动的瞬心；若其中一个刚体固定，则接触点即为运动刚度的绝对瞬心。
\end{itemize}



\paragraph{方法四：速度影像法 (Velocity Image Method) 的逆向应用}
根据公式 $\vect{v}_i = \vect{\omega} \times \vect{r}_{Pi}$，刚体上各点速度的大小与其到瞬心的距离成正比。
\begin{equation}
    \frac{v_A}{PA} = \frac{v_B}{PB} = \frac{v_C}{PC} = \omega
\end{equation}
通过建立比例关系方程组，可以解析地解出瞬心坐标 $(x_P, y_P)$。

\paragraph{瞬心位置的特殊情况汇总}
\begin{longtable}{ll}
    \toprule
    \textbf{运动状态} & \textbf{瞬心位置 $P$} \\ \midrule
    定轴转动 & 位于转轴中心（固定不变） \\
    瞬时平动 ($\omega=0$) & 位于无穷远处 \\
    纯滚动圆盘 & 位于与地面的接触点 \\
    滑块沿固定平面滑动 & 位于过滑块接触面法线的无穷远处 \\ \bottomrule
\end{longtable}

\subsection{常见错误辨析：速度瞬心 vs 加速度瞬心}
\begin{tcolorbox}[colback=red!5!white, colframe=red!75!black, title=绝对禁止的操作]
    **错误辨析**：很多学生认为速度瞬心的加速度也为零。\\
    **正确物理图像**：速度瞬心 $P$ 只是在这一瞬时速度为零。由于刚体仍在转动，该点通常具有向心加速度。加速度为零的点称为“加速度瞬心”，其位置与速度瞬心完全不同。在计算加速度时，严禁将速度瞬心当做定轴转动的轴心使用！
\end{tcolorbox}
\subsection{加速度分析：基点法}
加速度不具有“瞬心”性质（瞬时加速度为零的点与速度瞬心不重合）。
\begin{equation}
    \vect{a}_A = \vect{a}_B + \vect{a}_{AB}^n + \vect{a}_{AB}^t
\end{equation}
其中：
\begin{itemize}
    \item $\vect{a}_{AB}^n = \vect{\omega} \times (\vect{\omega} \times \vect{r}_{BA})$，大小为 $\omega^2 \cdot AB$，指向基点 $B$。
    \item $\vect{a}_{AB}^t = \vect{\alpha} \times \vect{r}_{BA}$，大小为 $\alpha \cdot AB$，垂直于 $AB$ 连线。
\end{itemize}

\subsection{经典例题分析：曲柄连杆机构}
设曲柄 $OA$ 以等角速度 $\omega_0$ 转动，连杆 $AB$ 长度为 $L$。求连杆中点 $M$ 的速度与加速度。
\begin{enumerate}
    \item \textbf{速度分析}：
    利用基点法或瞬心法确定连杆 $AB$ 的角速度 $\omega_{AB}$。
    $\vect{v}_A = \vect{v}_O + \vect{\omega}_0 \times \vect{r}_A$。
    \item \textbf{加速度分析}：
    以 $A$ 为基点分析 $B$ 点加速度，建立方程：
    $\vect{a}_B = \vect{a}_A + \vect{a}_{BA}^n + \vect{a}_{BA}^t$。
    通过几何投影解出连杆角加速度 $\alpha_{AB}$。
\end{enumerate}

\section{运动学常见辨析与误区}
\begin{tcolorbox}[colback=orange!5!white, colframe=orange!75!black, title=辨析：牵连加速度的误区]
    牵连加速度 $\vect{a}_e$ **不是**动系原点的加速度。\\
    它是“牵连点”在定系中的加速度。若动系绕定点转动，牵连点作圆周运动，故 $\vect{a}_e$ 包含向心项和切向项：$\vect{a}_e = \vect{a}_O' + \vect{\alpha}_e \times \vect{r}_r + \vect{\omega}_e \times (\vect{\omega}_e \times \vect{r}_r)$。
\end{tcolorbox}


\section{动力学基本定律 (Basic Laws)}

动力学研究受力、质量与运动之间的因果关系。其基础是牛顿三定律在矢量空间中的表达。

\subsection{质点动力学微分方程}
根据牛顿第二定律，质点在惯性参考系下的运动方程为：
\begin{equation}
    m \vect{a} = \sum \vect{F}_i
\end{equation}
在不同坐标系下的投影形式：
\begin{itemize}
    \item \textbf{直角坐标系}：$m\ddot{x} = \sum F_x, \quad m\ddot{y} = \sum F_y, \quad m\ddot{z} = \sum F_z$
    \item \textbf{自然坐标系}：
    \begin{align}
        m \der{v}{t} &= \sum F_{\tau} \quad (\text{切向：改变速度大小}) \\
        m \frac{v^2}{\rho} &= \sum F_n \quad (\text{法向：改变速度方向})
    \end{align}
\end{itemize}

% ==============================================================================
% 第二章：质点系动力学普遍定理 (The General Theorems)
% ==============================================================================
\section{质点系动力学普遍定理}

这是矢量力学的核心，通过对质量系统进行积分，得到描述整体运动规律的三个物理量：动量、动量矩、动能。

\subsection{动量定理 (Theorem of Momentum)}
\textbf{定义}：质点系的动量 $\vect{P} = \sum m_i \vect{v}_i = M \vect{v}_C$（其中 $C$ 为质心）。

\textbf{微分形式}：
\begin{equation}
    \der{\vect{P}}{t} = \sum \vect{F}_e
\end{equation}
\textbf{质心运动定理}：$M \vect{a}_C = \sum \vect{F}_e$。
\textit{推论}：若作用于系统的外力矢量和为零，则质心保持静止或匀速直线运动。

\subsection{动量矩定理 (Theorem of Angular Momentum)}
\textbf{定义}：质点系对定点 $O$ 的动量矩 $\vect{L}_O = \sum (\vect{r}_i \times m_i \vect{v}_i)$。
\textbf{基本方程}：
\begin{equation}
    \der{\vect{L}_O}{t} = \sum \vect{M}_O(\vect{F}_e)
\end{equation}

\subsubsection{刚体定轴转动与平面运动的表达}
\begin{itemize}
    \item \textbf{定轴转动}：$L_z = J_z \omega \Rightarrow J_z \alpha = \sum M_z(\vect{F}_e)$。
    \item \textbf{平面运动}：对质心 $C$ 列方程，$\vect{L}_C = J_C \vect{\omega}$。
\end{itemize}

\subsection{动能定理 (Theorem of Kinetic Energy)}
\textbf{定义}：系统的动能 $T = \sum \frac{1}{2} m_i v_i^2$。
\textbf{柯尼希定理 (Koenig's Theorem)}：
\begin{equation}
    T = \frac{1}{2} M v_C^2 + \frac{1}{2} J_C \omega^2
\end{equation}
即：平动动能 + 绕质心转动的转动动能。

\textbf{积分形式}：$T_2 - T_1 = \sum W_e + \sum W_i$。
\textit{注：对于刚体，由于内部约束力不做功，内力功之和 $\sum W_i = 0$。}

% ==============================================================================
% 第三章：刚体几何性质与转动惯量
% ==============================================================================
\section{刚体几何性质：转动惯量深度剖析}

转动惯量是衡量刚体转动惯性的度量。

\subsection{平行轴定理 (Huygens-Steiner Theorem)}
若已知刚体对通过质心 $C$ 的轴 $z_C$ 的转动惯量为 $J_C$，则对与之平行的轴 $z$ 的转动惯量为：
\begin{equation}
    J_z = J_C + Md^2
\end{equation}
其中 $d$ 为两轴间的距离。

\subsection{常见形状转动惯量表}
\begin{longtable}{lll}
    \toprule
    \textbf{刚体形状} & \textbf{轴线位置} & \textbf{转动惯量 $J$} \\ \midrule
    细长杆 ($M, L$) & 通过质心且垂直杆轴 & $\frac{1}{12}ML^2$ \\
    均质圆盘 ($M, R$) & 通过质心且垂直盘面 & $\frac{1}{2}MR^2$ \\
    均质球体 ($M, R$) & 通过球心 & $\frac{2}{5}MR^2$ \\
    空心薄圆环 ($M, R$) & 通过中心垂直平面 & $MR^2$ \\ \bottomrule
\end{longtable}

% ==============================================================================
% 第四章：达朗贝尔原理 (动静法)
% ==============================================================================
\section{达朗贝尔原理：动静法 (The Principle of D'Alembert)}

这是将动力学问题转化为静力学平衡方程的神奇方法。

\subsection{惯性力及其简化}
\textbf{质点惯性力}：$\vect{F}_I = -m \vect{a}$。
\textbf{刚体惯性力系的简化结果（向质心 $C$ 简化）}：
\begin{enumerate}
    \item \textbf{惯性力主矢}：$\vect{\Phi}^R = -M \vect{a}_C$。
    \item \textbf{惯性力主矩}：$\vect{M}_C^I = -J_C \vect{\alpha}$。
\end{enumerate}

\subsection{平衡方程形式}
在动力学过程中，施加在刚体上的外力与惯性力构成平衡力系：
\begin{equation}
    \begin{cases}
        \sum \vect{F}_e + \vect{\Phi}^R = 0 \\
        \sum \vect{M}_O(\vect{F}_e) + \vect{M}_O^I = 0
    \end{cases}
\end{equation}

% ==============================================================================
% 第五章：经典题型分析方法
% ==============================================================================
\subsection{经典题型分析：带摩擦的纯滚动问题}

这是哈工大最典型的综合考题：一个圆柱体在斜面上受力滚动，判断其运动状态并求加速度。

\subsubsection{解题逻辑流程}
\begin{enumerate}
    \item \textbf{假设状态}：首先假设圆柱体为“纯滚动”。此时满足约束方程 $a_C = \alpha R$。
    \item \textbf{列动力学方程}：
    \begin{align}
        M a_C &= Mg \sin \theta - F_f \\
        J_C \alpha &= F_f \cdot R
    \end{align}
    \item \textbf{联立求解}：解出加速度 $a_C$ 和所需的静摩擦力 $F_f$。
    \item \textbf{校验假设}：检查是否满足 $F_f \leq f_s F_n$。若满足，假设成立；若不满足，则刚体作“带滑动的滚动”，此时 $F_f = \mu F_n$，且 $a_C \neq \alpha R$。
\end{enumerate}

% ==============================================================================
% 第六章：常见误区辨析
% ==============================================================================
\section{动力学核心辨析：避坑指南}

\begin{tcolorbox}[colback=orange!5!white, colframe=orange!75!black, title=辨析 01：内力做功问题]
    虽然刚体的内力功之和为零，但**质点系**的内力功不一定为零。例如，两个用弹簧连接的滑块在运动时，弹簧的内力（弹力）会对系统做功。只有当质点间距离保持不变（即刚体模型）时，内力功才消失。
\end{tcolorbox}

\begin{tcolorbox}[colback=red!5!white, colframe=red!75!black, title=辨析 02：瞬心的加速度]
    在动力学中，速度瞬心 $P$ 的速度为零，但其加速度 $\vect{a}_P$ 通常不为零。因此，**不能**像静力学那样对速度瞬心列取矩方程 $\sum M_P = J_P \alpha$（除非该点是固定点或质心）。正确的取矩方程矩心应优先选择质心或固定转轴。
\end{tcolorbox}

\section{单自由度系统的线性振动}

机械振动是指物体在平衡位置附近的往复运动。矢量力学的核心任务是建立运动微分方程。

\subsection{自由振动 (Free Vibration)}
考虑一个由质量 $m$ 和刚度 $k$ 组成的系统，忽略阻力。取平衡位置为坐标原点 $x=0$。

\subsubsection{运动方程的建立 (利用牛顿第二定律)}
根据达朗贝尔原理或牛顿第二定律：
\begin{equation}
    m\dder{x}{t} = -kx \implies m\ddot{x} + kx = 0
\end{equation}
令 $\omega_n = \sqrt{k/m}$ 为系统的**固有角频率**，方程变为标准型：
\begin{equation}
    \ddot{x} + \omega_n^2 x = 0
\end{equation}

\subsubsection{运动规律}
该方程的通解为：
\begin{equation}
    x(t) = A \cos(\omega_n t + \phi)
\end{equation}
其中，$A$ 为振幅，$\phi$ 为初相位，均由初始条件 $x_0$ 和 $v_0$ 决定。



\subsection{衰减振动 (Damped Vibration)}
考虑线性粘性阻尼力 $F_d = -c\dot{x}$，其中 $c$ 为阻尼系数。
\textbf{运动方程}：
\begin{equation}
    m\ddot{x} + c\dot{x} + kx = 0 \implies \ddot{x} + 2n\dot{x} + \omega_n^2 x = 0
\end{equation}
其中 $n = c/2m$ 为衰减系数。根据 $n$ 与 $\omega_n$ 的关系，分为三种状态：
\begin{itemize}
    \item \textbf{欠阻尼 ($n < \omega_n$)}：系统作振幅不断减小的对数衰减振动。
    \item \textbf{过阻尼 ($n > \omega_n$)}：系统不发生振动，缓慢回到平衡位置。
    \item \textbf{临界阻尼 ($n = \omega_n$)}：系统最快回到平衡位置且不超越。
\end{itemize}

\subsection{受迫振动 (Forced Vibration)}
系统在简谐激振力 $F(t) = F_0 \sin(\omega t)$ 作用下的运动。
\textbf{全解结构}：$x(t) = x_h(t) + x_p(t)$（伴随振动 + 稳态振动）。
稳态振动的振幅 $X$ 为：
\begin{equation}
    X = \frac{F_0/k}{\sqrt{(1 - \lambda^2)^2 + (2\zeta\lambda)^2}}
\end{equation}
其中 $\lambda = \omega/\omega_n$ 为频率比，$\zeta = n/\omega_n$ 为阻尼比。

\begin{tcolorbox}[colback=red!5, colframe=red!75!black, title=共振 (Resonance)]
    当激振频率 $\omega$ 接近固有频率 $\omega_n$（即 $\lambda \to 1$）时，振幅急剧增大。在无阻尼情况下，理论振幅趋于无穷大。这是工程设计中必须规避或利用的状态。
\end{tcolorbox}



% ==============================================================================
% 第二部分：碰撞 (Impact)
% ==============================================================================
\section{碰撞专题 (Theory of Impact)}

碰撞的特征：时间极短 ($\Delta t \to 0$)，相互作用力（碰撞力）极大，导致速度发生有限的突变。

\subsection{碰撞的基本假设}
\begin{enumerate}
    \item \textbf{忽略微小力}：在碰撞瞬间，有限大小的力（如重力、普通弹簧力）产生的冲量忽略不计。
    \item \textbf{位移不变}：碰撞瞬间，物体的位置和构形保持不变。
\end{enumerate}

\subsection{恢复系数 (Coefficient of Restitution)}
碰撞过程分为：\textbf{变形阶段}（速度趋同）和\textbf{恢复阶段}（速度分离）。
恢复系数 $e$ 定义为分离速度与接近速度的比值：
\begin{equation}
    e = \frac{v_{2\tau} - v_{1\tau}}{u_{1\tau} - u_{2\tau}}
\end{equation}
其中 $u$ 为碰前速度，$v$ 为碰后速度。
\begin{itemize}
    \item $e = 1$：完全弹性碰撞（动能无损失）。
    \item $e = 0$：完全非弹性碰撞（碰撞后两物体粘在一起）。
    \item $0 < e < 1$：典型塑性碰撞。
\end{itemize}



\subsection{碰撞动力学方程}
\subsubsection{质点系冲量定理}
\begin{equation}
    M(\vect{v}_{C2} - \vect{v}_{C1}) = \sum \vect{S}_e
\end{equation}
其中 $\vect{S}_e$ 为外部碰撞冲量。

\subsubsection{对某定点 $O$ 的冲量矩定理}
\begin{equation}
    \vect{L}_{O2} - \vect{L}_{O1} = \sum \vect{M}_O(\vect{S}_e)
\end{equation}
对于定轴转动刚体：$J_z(\omega_2 - \omega_1) = \sum M_z(\vect{S}_e)$。

\subsection{ 撞击中心 (Center of Percussion)}
当刚体受碰撞冲量作用时，若能使转轴处的碰撞约束反力为零，则该冲量的作用点称为撞击中心。
设转轴到质心的距离为 $l_C$，则撞击中心到转轴的距离 $l_P$ 满足：
\begin{equation}
    l_P = \frac{J_z}{M l_C}
\end{equation}
\textit{工程应用：锤子、棒球棒的“甜蜜点”即利用此原理减少手部震动。}

% ==============================================================================
% 第三部分：综合解析
% ==============================================================================
\subsection{常见辨析与解题范式}

\begin{tcolorbox}[colback=blue!5, colframe=blue!75!black, title=解题策略：振动 vs 碰撞]
    \textbf{振动题目}：关键在于“找回复力”。推荐使用动能定理对时间求导法（能量法）建立微分方程，可自动消除不作功的约束力。\\
    \textbf{碰撞题目}：关键在于“分清突变”。必须画出碰前、碰后的速度矢量图，并应用冲量矩定理，注意恢复系数方程是必不可少的补充方程。
\end{tcolorbox}

\end{document}